% Experiment report. Fixed sections: see experiments/TEMPLATE/report.tex and experiments/README.md.
% Build: npm run experiment:report -- 04-road-aware-moves
\documentclass{../polyexp}
\begin{document}
% generated by scripts/experiment-report.ts from experiment.json and models/<run>/ - do not edit
\providecommand{\ExpHeader}{}\renewcommand{\ExpHeader}{%
{\sffamily\footnotesize\color{expgrey}EXPERIMENT 4 \textperiodcentered\ MECHANISM \hfill 2026-11-24}\par
{\LARGE\bfseries\sffamily Road-aware moves}\hfill\verdictbox{expred}{NOT KEPT}\par\smallskip
{\itshape Does the bot get stronger when moves count real movement cost, so that roads are used?}\par\medskip
\begin{tabular}{@{}p{22mm}p{\dimexpr\linewidth-26mm}@{}}
\fact{Change}{A move's progress becomes tiles gained + roads $\times$ (road-aware gain $-$ tiles gained); steps between two road tiles cost half, as in the engine.}
\fact{Knobs}{\knob{roads} $\in$ \{0.25, 0.5, 1, 2\}, best 0.5}
\fact{Base}{\model{greedy-full-front}}
\fact{Method}{Strength sweep, 200 games per point against each panel model, final test 400 games; metric: points share}
\fact{Setups}{4 players, 16$\times$16, drylands, turn cap 150}
\fact{Record}{issue \issue{34}, pull request \issue{125}; status: done}
\end{tabular}\par\medskip
\begin{center}\begin{tikzpicture}[node distance=5mm]
    \node[flowbox] (n0) {{\sffamily\bfseries\scriptsize\color{expgrey}BASE}\\[1pt]\model{full-front}};
    \node[flowbox,right=5mm of n0] (n1) {{\sffamily\bfseries\scriptsize\color{expgrey}CHANGE}\\[1pt]\knob{roads}\\{\scriptsize 0.25, 0.5, 1, 2}};
    \node[flowbox,right=5mm of n1] (n2) {{\sffamily\bfseries\scriptsize\color{expgrey}MEASURED}\\[1pt]Strength sweep, 4 values\\{\scriptsize 1 setup, metric: points share}};
    \node[flowbox,right=5mm of n2] (n3) {{\sffamily\bfseries\scriptsize\color{expgrey}RESULT}\\[1pt]Points vs base, final test\\{\scriptsize 50.0 $\rightarrow$ 50.2\,\%}\\\textbf{+0.2\pp}};
    \node[flowbox,right=5mm of n3] (n4) {{\sffamily\bfseries\scriptsize\color{expgrey}VERDICT}\\[1pt]\textbf{\color{expred}NOT KEPT}};
    \draw[flowarrow] (n0) -- (n1); \draw[flowarrow] (n1) -- (n2); \draw[flowarrow] (n2) -- (n3); \draw[flowarrow] (n3) -- (n4);
\end{tikzpicture}\end{center}
\begin{center}\sffamily\small\begin{tabular}{@{}lrrr@{}}\toprule
Measure & Base & New & Change\\\midrule
Points vs base, final test & 50.0\,\% & 50.2\,\% & +0.2\\
Won vs base, final test & 50.0\,\% & 50.7\,\% & +0.7\\
Points share, 7-model panel & 67.4\,\% & 67.4\,\% & +0.0\\
\bottomrule\end{tabular}\end{center}
}
\providecommand{\ExpResults}{}\renewcommand{\ExpResults}{%
\subsection{Strength sweep {\normalfont\small(\model{greedy-full-roads})}}
\begin{center}\sffamily\small\begin{tabular}{@{}lrrrrrr@{}}\toprule
Point & Finish (\%) & Points (\%) & Won (\%) & Conquest (\%) & Turns & vs base: points (\%)\\\midrule
base & -- & 67.36 & 80.7 & -- & -- & 50.0\\
roads=0.25 & -- & 67.37 & 80.9 & -- & -- & 49.9\\
roads=0.5 & -- & 67.40 & 80.7 & -- & -- & 50.4\\
roads=1 & -- & 67.13 & 80.1 & -- & -- & 50.2\\
roads=2 & -- & 66.74 & 79.9 & -- & -- & 50.0\\
\bottomrule\end{tabular}\\[3pt]{\footnotesize\color{expgrey}Panel means; 200 games against each of 7 panel models per point.}\end{center}
\begin{center}\begin{tikzpicture}\begin{axis}[expplot,height=4.6cm,xtick={0,1,2,3,4},xticklabels={{base},{0.25},{0.5},{1},{2}},
  xlabel={\texttt{roads}},ylabel={Points share, panel mean (\%)},enlarge x limits=0.12,enlarge y limits=0.25]
\addplot[dashed,expgrey,line width=0.7pt,forget plot] coordinates {(0,67.357) (4,67.357)};
\addplot[only marks,mark=*,mark size=2.6pt,expblue] coordinates { (1,67.371) (2,67.400) (3,67.129) (4,66.743)};
\addplot[only marks,mark=*,mark size=2.6pt,expgrey] coordinates {(0,67.357)};
\end{axis}\end{tikzpicture}\\[-2pt]{\footnotesize\color{expgrey}Dashed: the base. Grey: base; blue: sweep points.}\end{center}
\begin{center}\sffamily\small\begin{tabular}{@{}lrr@{}}\toprule
Opponent & Won (\%) & Points (\%)\\\midrule
\model{v1} & 93 $\pm$ 2 & 75.5\\
\model{hill-v1} & 91 $\pm$ 3 & 72.3\\
\model{hill-v2} & 94 $\pm$ 2 & 73.1\\
\model{hill-v3} & 90 $\pm$ 3 & 69.7\\
\model{hill-v4} & 86 $\pm$ 3 & 68.6\\
\model{full-v1} & 57 $\pm$ 5 & 54.4\\
\model{full-front} & 51 $\pm$ 5 & 50.2\\
\midrule
\textbf{Mean without \model{full-front}} & 85.1 & 68.9\\
\bottomrule\end{tabular}\\[3pt]{\footnotesize\color{expgrey}Final test of roads=0.5: 400 unseen games against each.}\end{center}
}

\ExpHeader

\section{Question and hypothesis}
In a live 9-bot game at turn 40 the bots had built 166 road tiles (41\,\% of the map, about 330 stars), yet only 17 of 825 moves started and ended on the road network. Move scoring counts tiles gained and ignores movement cost. \textbf{Hypothesis:} scoring moves by road-aware distance makes units use roads and reach goals sooner.

\section{Mechanism}
A Dijkstra map from the unit's goal gives a road-aware distance in half-steps. Progress of a move $= t + \knob{roads}\cdot(r - t)$, with $t$ the tiles gained and $r$ the road-aware gain. A step between two road tiles (road, bridge, village or city on friendly or unclaimed land) costs half. Default 0. Code: \texttt{road\_dist} in \texttt{bot\_full.rs}.

\section{Setup}
Sanity check on the live game at \knob{roads} $=1$: 13\,\% of units pick a different best move, and slightly more moves end on a road or city (992 against 935 of 1,247). The panel has seven models here, because the base is added to the fixed six.

\section{Results}
\ExpResults

No effect at weak values, a small cost at strong ones (\knob{roads} 2: $-0.6$ points). Final test of the best point against the base: 50.2\,\% of points, 51\,\% won ($\pm$5).

\section{Behaviour}
Moves ending on a road or city at \knob{roads} $=1$: 992 of 1,247, against 935 for the base (one live position).

\section{Interpretation}
Roads change which tile a unit ends on, but rarely what it achieves. On 16$\times$16 maps most moves reach their goal within one or two turns, so a half-cost path saves little. The hypothesis that unused roads are a strength problem is not supported; they are a waste of stars, and stars are not the bottleneck (2,500 banked by the late game).

\todo{jakob}{Your reading: does this match what you saw? Anything to add or correct?}

\section{Threats to validity}
\begin{itemize}
\item Only 16$\times$16 maps; on larger maps paths are longer and roads could matter.
\item The other half of the idea (build roads only where they connect something) is untested.
\end{itemize}

\section{Verdict and next step}
\textbf{Not kept.} The code stays, switched off. The champion remains \model{greedy-full-front}. Next: a yardstick that rewards finishing (experiment 5).

\section{Notes for the paper}
Draft: a clean negative result. Correcting a visibly ``wrong'' behaviour (roads built but unused) does not change strength, because it is not on the critical path to winning.

\todo{jakob}{Is this worth a place in the paper? Where?}

\end{document}
